\documentclass[aspectratio=169]{beamer}

% ── Theme ────────────────────────────────────────────────────────────────────
\usetheme{Madrid}
\usecolortheme{seahorse}
\usefonttheme{professionalfonts}

% ── Packages ─────────────────────────────────────────────────────────────────
\usepackage{amsmath, amssymb}
\usepackage{booktabs}
\usepackage{xcolor}
\usepackage{tcolorbox}
\tcbuselibrary{skins, breakable}

% ── Colours ──────────────────────────────────────────────────────────────────
\definecolor{oxblue}{HTML}{002147}
\definecolor{noteblue}{HTML}{d0e8f5}
\definecolor{noteframe}{HTML}{045a8d}
\definecolor{warnbg}{HTML}{fff3cd}
\definecolor{warnframe}{HTML}{e6a817}
\definecolor{tipbg}{HTML}{d4edda}
\definecolor{tipframe}{HTML}{28a745}

\setbeamercolor{frametitle}{bg=oxblue, fg=white}
\setbeamercolor{title}{fg=oxblue}
\setbeamercolor{subtitle}{fg=oxblue!70}
\setbeamercolor{author}{fg=oxblue}
\setbeamercolor{institute}{fg=oxblue!80}
\setbeamercolor{date}{fg=oxblue!70}
\setbeamercolor{block title}{bg=oxblue!15, fg=oxblue}
\setbeamercolor{block body}{bg=oxblue!5}

% ── Callout boxes ────────────────────────────────────────────────────────────
\newenvironment{callout-note}[1][Note]{%
  \begin{tcolorbox}[colback=noteblue, colframe=noteframe,
    title={\textbf{#1}}, fonttitle=\small\bfseries,
    left=4pt, right=4pt, top=2pt, bottom=2pt]
}{\end{tcolorbox}}

\newenvironment{callout-warning}[1][Warning]{%
  \begin{tcolorbox}[colback=warnbg, colframe=warnframe,
    title={\textbf{#1}}, fonttitle=\small\bfseries,
    left=4pt, right=4pt, top=2pt, bottom=2pt]
}{\end{tcolorbox}}

\newenvironment{callout-tip}[1][Tip]{%
  \begin{tcolorbox}[colback=tipbg, colframe=tipframe,
    title={\textbf{#1}}, fonttitle=\small\bfseries,
    left=4pt, right=4pt, top=2pt, bottom=2pt]
}{\end{tcolorbox}}

% ── Title ─────────────────────────────────────────────────────────────────────
\title{Panel Data Fundamentals}
\subtitle{Introduction to Longitudinal Data Analysis with R}
\author{Dr Clemens Jarnach}
\institute{University of Oxford}
\date{2025--26}

% ═════════════════════════════════════════════════════════════════════════════
\begin{document}

\begin{frame}
  \titlepage
\end{frame}

% ── Slide 1 ───────────────────────────────────────────────────────────────────
\begin{frame}{What is Panel Data?}
  \begin{itemize}
    \item A dataset that tracks the \textbf{same units} (individuals, firms, countries) \textbf{over time}
    \item Also called \textbf{longitudinal data} or \textbf{cross-sectional time-series} data
    \item Combines two dimensions:
      \begin{itemize}
        \item \textbf{Cross-sectional}: variation across units $i$
        \item \textbf{Time-series}: variation over time $t$
      \end{itemize}
  \end{itemize}

  \vspace{0.5em}
  \begin{callout-note}
    Each observation is indexed by both a unit $i$ and a time period $t$: $(y_{it}, x_{it})$
  \end{callout-note}
\end{frame}

% ── Slide 2 ───────────────────────────────────────────────────────────────────
\begin{frame}{Why Panel Data? Key Advantages}
  \begin{itemize}
    \item Combines \textbf{cross-sectional + time-series variation}
    \item Enables:
      \begin{itemize}
        \item<2-> Control for \textbf{individual heterogeneity}
        \item<3-> More \textbf{efficient estimation} (greater variability)
        \item<4-> Analysis of \textbf{dynamic processes}
        \item<5-> Identification of effects \textbf{not visible in cross-sections}
        \item<6-> Use of \textbf{micro-level data} (better measurement)
        \item<7-> Cross-dimensional inference (time $\leftrightarrow$ individuals)
      \end{itemize}
  \end{itemize}
\end{frame}

% ── Slide 3 ───────────────────────────────────────────────────────────────────
\begin{frame}{The Core Problem: Unobserved Heterogeneity}
  Consider the model:
  \[
    y_{it} = \alpha + \beta x_{it} + \gamma z_i + u_{it}
  \]
  \begin{itemize}
    \item $z_i$ is an \textbf{unobserved}, time-invariant characteristic
    \item If $z_i$ is correlated with $x_{it}$ $\Rightarrow$ \textbf{omitted variable bias}
    \item OLS estimates of $\beta$ are \textbf{inconsistent}
  \end{itemize}

  \vspace{0.5em}
  \begin{callout-warning}
    \textbf{Key condition for consistency:} $z_i$ must be uncorrelated with regressors and the error term.
  \end{callout-warning}
\end{frame}

% ── Slide 4 ───────────────────────────────────────────────────────────────────
\begin{frame}{Intuition: Agricultural Example}
  Suppose farm output depends on:

  \vspace{0.5em}
  \begin{center}
    \begin{tabular}{ll}
      \toprule
      \textbf{Variable} & \textbf{Status} \\
      \midrule
      Labour $x_{it}$ & Observed \\
      Soil quality $z_i$ & Unobserved \\
      \bottomrule
    \end{tabular}
  \end{center}

  \vspace{0.8em}
  \textbf{The problem:}
  \begin{itemize}
    \item Farmers likely adjust labour based on their soil quality
    \item $x_{it}$ is correlated with $z_i$
    \item $\Rightarrow$ \textbf{OLS becomes inconsistent}
  \end{itemize}
\end{frame}

% ── Slide 5 ───────────────────────────────────────────────────────────────────
\begin{frame}{The General Panel Data Model}
  Absorb unobserved heterogeneity into an individual-specific intercept:
  \[
    y_{it} = \alpha_i + \beta x_{it} + u_{it}
  \]
  \begin{itemize}
    \item $\alpha_i$ captures all \textbf{time-invariant unobserved factors} for unit $i$
    \item The goal: \textbf{remove or control for} $\alpha_i$
  \end{itemize}

  \vspace{0.5em}
  \onslide<2->{%
    \begin{block}{}
      This is the starting point for both \textbf{Fixed Effects} and \textbf{Random Effects} models.
    \end{block}
  }
\end{frame}

% ── Section divider ──────────────────────────────────────────────────────────
\begin{frame}[plain]
  \begin{center}
    \vspace{2.5em}
    {\Large\bfseries\color{oxblue} Three Strategies for Eliminating $\alpha_i$}
  \end{center}
\end{frame}

% ── Slide 6 ───────────────────────────────────────────────────────────────────
\begin{frame}{Strategy 1: First-Difference (FD) Estimator}
  \textbf{Idea:} Subtract last period's equation from the current period:
  \[
    \Delta y_{it} = \beta\,\Delta x_{it} + \Delta u_{it}
  \]
  where $\Delta y_{it} = y_{it} - y_{i,t-1}$

  \vspace{0.8em}
  \begin{columns}[T]
    \begin{column}{0.5\textwidth}
      \textbf{Advantages}
      \begin{itemize}
        \item Eliminates $\alpha_i$ completely
        \item Estimable via pooled OLS
      \end{itemize}
    \end{column}
    \begin{column}{0.5\textwidth}
      \textbf{Disadvantages}
      \begin{itemize}
        \item Loses first time period
        \item Can amplify measurement error
      \end{itemize}
    \end{column}
  \end{columns}
\end{frame}

% ── Slide 7 ───────────────────────────────────────────────────────────────────
\begin{frame}{Strategy 2: Least Squares Dummy Variables (LSDV)}
  \textbf{Approach:} Include a dummy variable $d_i$ for each unit:
  \[
    y_{it} = \sum_i \alpha_i d_i + \beta x_{it} + u_{it}
  \]
  \begin{itemize}
    \item<2-> Estimates a \textbf{separate intercept} $\alpha_i$ for each unit
    \item<3-> $\hat{\beta}$ is $\sqrt{NT}$-consistent
    \item<4-> Individual effects estimated with $\sqrt{T}$-consistency
    \item<5-> \textbf{Drawback:} $N$ additional parameters $\rightarrow$ loss of degrees of freedom; computationally costly for large $N$
  \end{itemize}
\end{frame}

% ── Slide 8 ───────────────────────────────────────────────────────────────────
\begin{frame}{Strategy 3: Fixed Effects (Within) Estimator}
  \textbf{Idea:} Remove $\alpha_i$ by \textbf{demeaning} each variable over time:
  \[
    \tilde{y}_{it} = y_{it} - \bar{y}_{i}, \qquad \tilde{x}_{it} = x_{it} - \bar{x}_{i}
  \]
  This yields the \textbf{within-transformed} model:
  \[
    \tilde{y}_{it} = \beta\,\tilde{x}_{it} + \tilde{u}_{it}
  \]

  \vspace{0.4em}
  \begin{callout-tip}
    $\alpha_i$ cancels out because $\bar{\alpha}_i = \alpha_i$. Only \textbf{within-unit variation over time} identifies $\beta$.
  \end{callout-tip}
\end{frame}

% ── Slide 9 ───────────────────────────────────────────────────────────────────
\begin{frame}{Fixed Effects: Properties and Limitations}
  \begin{columns}[T]
    \begin{column}{0.55\textwidth}
      \textbf{Advantages}
      \begin{itemize}
        \item Eliminates $\alpha_i$ without dummy variables
        \item Computationally efficient
        \item Consistent even when $\alpha_i$ is correlated with $x_{it}$
        \item Standard approach in applied work
      \end{itemize}
    \end{column}
    \begin{column}{0.45\textwidth}
      \textbf{Limitations}
      \begin{itemize}
        \item Cannot estimate effects of \textbf{time-invariant} variables (e.g.\ gender, country)
        \item Relies solely on \textbf{within-unit} variation
        \item Less efficient when $\alpha_i$ is uncorrelated with $x_{it}$
      \end{itemize}
    \end{column}
  \end{columns}
\end{frame}

% ── Slide 10 ──────────────────────────────────────────────────────────────────
\begin{frame}{FD vs LSDV vs Within Estimator}
  \begin{center}
    \begin{tabular}{llll}
      \toprule
      \textbf{Method} & \textbf{Removes $\alpha_i$?} & \textbf{Parameters} & \textbf{Notes} \\
      \midrule
      First-difference & Yes & None added  & Loses one period \\
      LSDV             & Yes & $N$ dummies & Costly for large $N$ \\
      Within (FE)      & Yes & None added  & Standard choice \\
      \bottomrule
    \end{tabular}
  \end{center}

  \vspace{0.6em}
  \onslide<2->{%
    \begin{block}{}
      All three are \textbf{algebraically equivalent} under strict exogeneity ---
      the within estimator is preferred for computational efficiency.
    \end{block}
  }
\end{frame}

% ── Slide 11 ──────────────────────────────────────────────────────────────────
\begin{frame}{Key Takeaways}
  \begin{enumerate}
    \item Panel data solves \textbf{unobserved heterogeneity} by exploiting \textbf{within-unit variation over time}
    \item The core challenge: unobserved $\alpha_i$ correlated with regressors $\Rightarrow$ OLS bias
    \item Three equivalent strategies: \textbf{first-differencing}, \textbf{LSDV}, \textbf{within (demeaning)}
    \item \textbf{Fixed Effects = most widely used practical method}
  \end{enumerate}

  \vspace{0.5em}
  \begin{callout-note}
    Next: when should we use \textbf{Random Effects} instead, and how do we choose?
  \end{callout-note}
\end{frame}

\end{document}
